\subsection{\texorpdfstring{同态 \(\operatorname{Tor}^{\mathbf{B}}(\mathbf{P},\mathbf{N}\otimes_{\mathbf{A}}\mathbf{Q})\to \operatorname{Tor}^{\mathbf{A}}(\mathbf{P}\otimes_{\mathbf{B}}\mathbf{N},\mathbf{Q})\) 与 \(\operatorname{Ext}_{\mathbf{A}}(\mathbf{Q},\operatorname{Hom}_{\mathbf{B}}(\mathbf{N},\mathbf{M}))\to \operatorname{Ext}_{\mathbf{B}}(\mathbf{N}\otimes_{\mathbf{A}}\mathbf{Q},\mathbf{M})\)}{Tor 与 Ext 的换环同态}}

沿用前面的记号，并假设 N 在 A 上是平坦的。则态射 \(\mathrm{N}\otimes_{\mathbf{A}}\mathrm{L}(\mathrm{Q})\xrightarrow{\mathrm{1}\otimes p_{\mathrm{Q}}} \mathrm{N}\otimes_{\mathbf{A}}\mathrm{Q}\) 是一个拟同构（X，第 67 页，命题 4），因此同态

\[
\psi (\mathrm{P}, \mathrm{N} \otimes_ {\mathrm{A}} \mathrm{L} (\mathrm{Q})): \operatorname{Tor} ^ {\mathrm{B}} (\mathrm{P}, \mathrm{N} \otimes_ {\mathrm{A}} \mathrm{Q}) \rightarrow \mathrm{H} (\mathrm{P} \otimes_ {\mathrm{B}} \mathrm{N} \otimes_ {\mathrm{A}} \mathrm{L} (\mathrm{Q}))
\]

\[
\varphi (\mathrm{N} \otimes_ {\mathrm{A}} \mathrm{L} (\mathrm{Q}), \mathrm{M}): \mathrm{H} (\operatorname{Homgr} _ {\mathrm{B}} (\mathrm{N} \otimes_ {\mathrm{A}} \mathrm{L} (\mathrm{Q}), \mathrm{M})) \rightarrow \operatorname{Ext} _ {\mathrm{B}} (\mathrm{N} \otimes_ {\mathrm{A}} \mathrm{Q}, \mathrm{M}).
\]

再利用同构

\[
\begin{array}{c} \overline {{\psi}} _ {Q} (P \otimes_ {B} N): \operatorname{Tor} ^ {A} (P \otimes_ {B} N, Q) \to H (P \otimes_ {B} N \otimes_ {A} L (Q)), \\ \beta : \operatorname{Homgr} _ {A} (L (Q), \operatorname{Hom} _ {B} (N, M)) \to \operatorname{Homgr} _ {B} (N \otimes_ {A} L (Q), M), \\ \varphi_ {Q} (\operatorname{Hom} _ {B} (N, M)): H (\operatorname{Homgr} _ {A} (L (Q), \operatorname{Hom} _ {B} (N, M))) \to \operatorname{Ext} _ {A} (Q, \operatorname{Hom} _ {B} (N, M)), \end{array}
\]

由此，我们推出次数为 0 的典范分次同态：

\begin{enumerate}
\def\labelenumi{(\arabic{enumi})}
\setcounter{enumi}{10}
\tightlist
\item
\end{enumerate}

\[
\operatorname{Tor} ^ {\mathbf {B}} (\mathrm{P}, \mathrm{N} \otimes_ {\mathrm{A}} \mathrm{Q}) \rightarrow \operatorname{Tor} ^ {\mathbf {A}} (\mathrm{P} \otimes_ {\mathrm{B}} \mathrm{N}, \mathrm{Q})\tag{12}
\]

\[
\operatorname{Ext} _ {\mathbf {A}} \left(\mathrm{Q}, \operatorname{Hom} _ {\mathbf {B}} (\mathrm{N}, \mathrm{M})\right)\rightarrow \operatorname{Ext} _ {\mathbf {B}} \left(\mathrm{N} \otimes_ {\mathbf {A}} \mathrm{Q}, \mathrm{M}\right).
\]

命题 8. --- a) 若 N 在 A 上和 B 上都是平坦的，则同态 (11) 是双射的。b) 若 N 在 A 上是平坦的且在 B 上是投射的，则同态 (12) 是双射的。

事实上， \(\mathbf{N} \otimes_{\mathbf{A}} \mathbf{L}(\mathbf{Q})\) 同构于若干份 \(\mathbf{N}\) 的直和，因此当 B-模 \(\mathbf{N}\) 是平坦的（相应地，投射的）时，它是平坦的（相应地，投射的）B-模；然后应用定理 1（ \(\mathbf{X}\) ，第 100 页）。

